\fancyfoot[C]{}
\fancypagestyle{plain}{\fancyfoot[C]{}}
\part{Rechtliches}
\label{part:legal}

\chapter{Rechtliche Einschränkungen}
\label{ch:legal_boundaries}
\fancyfoot[C]{\large Timofey Luzin}
\fancypagestyle{plain}{\fancyfoot[C]{\large Timofey Luzin}}

Da beim Projekt "`ALBERT"' Arbeit mit Pyrotechnik vorgesehen war, musste dringend beachtet werden, dass keine gesetzwidrigen Aktionen durchgeführt werden. Hierbei stoßte man auf das Problem, dass der Raketenmotor, der sich in der Motorklasse H befindet, in Österreich ohne einer speziellen Lizent weder gebaut, noch gezündet werden darf. Auch bei den pyrotechnischen Charges für den Fallschirmauswurf stoßte man auf rechtliche Barrieren.

\section{Raketenmotor}
\label{sec:rocket_motor_legal}
Modellraketenmotoren fallen in Österreich grundsätzlich in den Bereich des Pyrotechnikgesetzes. In der Arbeitsrichtlinie Pyrotechnische Gegenstände des Bundesfinanzministeriums für Finanzen findet man Modellbauraketenmotoren unter dem Punkt 1.1.3 "`Sonstige pyrotechnische Gegenstände"'.
\\
\\
"`Sonstige pyrotechnische Gegenstände sind alle pyrotechnischen Gegenstände, die keine
Feuerwerkskörper (Abschnitt 1.1.1.) und keine pyrotechnischen Gegenstände für Bühne und
Theater (Abschnitt 1.1.2.) sind. Sonstige pyrotechnische Gegenstände werden gemäß § 13
PyroTG 2010 entsprechend ihrer Verwendungsart oder ihrem Zweck und dem Grad ihrer
Gefährlichkeit einschließlich ihres Lärmpegels unterteilt in:
\begin{enumerate}
    \item \textbf{Kategorie P1}: Sonstige pyrotechnische Gegenstände, die eine geringe Gefahr darstellen. Von der Kategorie P1 sind zB Anzündschnüre, -lichter und -litzen, mechanische Anzünder, pyrotechnische Signalmittel (Berg- und Seenotsignalmittel, Signalstifte mit Munition, Munition von Leuchtpistolen, Rauchsignale, alpine Notsignalmittel, Hand-, Warn-, Bengal- und Magnesiumfackeln), Rauch- und Nebelerzeuger, pyrotechnische Gegenstände für die Schädlingsbekämpfung in der Land- und Forstwirtschaft, Gurtenstraffer, Airbags und Stromkreisunterbrecher erfasst, die als geringgefährlich gelten.
    \item \textbf{Kategorie P2}: Sonstige pyrotechnische Gegenstände, die Personen mit Fachkenntnissen vorbehalten sind. Der Kategorie P2 sind insbesondere Stoppinen, Anzündbänder („Tape- Match“), Blitzknallpatronen, Modellbauraketenmotoren, Hagel- und Starenabwehrraketen sowie Verzögerungsanzünder zuzuordnen, die eine Gefahr darstellen und daher nur von Personen mit Fachkenntnissen verwendet werden dürfen."'
\end{enumerate} \parencite{BMF_Pyro}
\newpage
Dieselbe Definition ist auch in auch in den Erläuterungen für das Pyrotechnikgesetz 2010 zu finden:
\\
"`\textbf{Zu § 22:} \\
Z 1: Von der Kategorie P1 sind zB Anzündschnüre, -lichter und -litzen, mechanische Anzünder, Gurtenstraffer, Airbags und Stromkreisunterbrecher erfasst, die als geringgefährlich gelten. \\
Z 2: Der Kategorie P2 sind insbesondere Stoppinen, Anzündbänder („Match-Tapes“), Blitzknallpatronen, Raketenmotoren und Verzögerungsanzünder zuzuordnen, die eine Gefahr darstellen."' \parencite{pyro_law_2010}
\\
\\
Anhand dieser Quellen ist klar zur erkennen, dass der für "`ALBERT"' vorgesehene Motor klar unter die Kategorie P2 fällt. Der Besitz und die Verwendung dieser Kategorie sind ohne einem Pyrotechnik-Ausweis und ohne einer behördlichen Bewilligung nicht erlaubt.
\\
"`(2) Besitz und Verwendung pyrotechnischer Gegenstände der Kategorie P2, mit Ausnahme pyrotechnischer Anzündmittel, sind nur aufgrund einer behördlichen Bewilligung erlaubt. Die Bewilligung ist nach Glaubhaftmachung eines Bedarfs für eine bestimmte Art (Produktgruppe) von pyrotechnischen Gegenständen der Kategorie P2 zu erteilen, wenn die Person gemäß Abs. 1 Z 2 lit. a oder b oder Z 3 über einen Pyrotechnik-Ausweis für die beantragte Art (Produktgruppe) der Kategorie P2 verfügt und unter Bedachtnahme auf die Umstände der beabsichtigten Verwendungen der pyrotechnischen Gegenstände, insbesondere die beantragten Verwendungsorte oder -gebiete und Verwendungszeiten gewährleistet ist, dass Gefährdungen von Leben, Gesundheit und Eigentum von Menschen oder der öffentlichen Sicherheit sowie unzumutbare Lärmbelästigungen vermieden werden. Die Berechtigung kann in Form einer Dauerbewilligung erteilt werden."' \parencite{pyroTG_jusline}
\\
\\
Da wir als Projektteam über keinen Pyrotechnik-Ausweis verfügen, darf der Raketenmotor weder gebaut, noch gezündet werden. Daher wurde das Projekt auf den Bau von der restlichen Rakete, auf die Entwicklung des Flightcomputers, die Strömungssimulation des Raketenkörpers und Entwicklung der Regelungstechnik beschränkt.

\section{Pyrotechnische Charges zum Fallschirmauswurf}
\label{sec:pyro_charges_legal}